\documentclass[12pt]{article}
\usepackage[american]{babel}
\usepackage[margin=1in]{geometry}
\usepackage{setspace}
\usepackage[T1]{fontenc}
\usepackage{newtxtext,newtxmath}
\usepackage{graphicx}
\usepackage{float}
\usepackage{booktabs}
\usepackage{tabularx}
\usepackage{array}
\usepackage{ragged2e}
\usepackage{caption}
\usepackage[table]{xcolor}
\usepackage{hyperref}
\usepackage{csquotes}
\usepackage[style=apa,backend=biber]{biblatex}
\usepackage{colortbl}
\usepackage{hhline}
\usepackage{tikz}
\usetikzlibrary{shapes.geometric, arrows.meta, positioning, fit, backgrounds}
\usepackage{xcolor}
\usepackage[T1]{fontenc}
\usepackage{lmodern}
\renewcommand{\familydefault}{\sfdefault}
\usepackage{tikz}
\usetikzlibrary{arrows.meta,positioning,calc}

\definecolor{collect}{HTML}{2F5E8C}
\definecolor{collectbg}{HTML}{EEF4FA}
\definecolor{clean}{HTML}{2F7A66}
\definecolor{cleanbg}{HTML}{EEF8F5}
\definecolor{prep}{HTML}{B87523}
\definecolor{prepbg}{HTML}{FCF3E8}
\definecolor{analysis}{HTML}{7A4F9E}
\definecolor{analysisbg}{HTML}{F5EFFA}
\definecolor{ink}{HTML}{1D1D1F}
\definecolor{muted}{HTML}{5C6166}
\definecolor{linegray}{HTML}{A8AEB6}
\definecolor{softgray}{HTML}{F5F6F7}

\tikzset{
  >={Latex[length=2.8mm,width=2.0mm]},
  stage/.style={rounded corners=7pt, line width=1.1pt, minimum height=13.2cm, inner sep=0pt},
  stagehead/.style={rounded corners=7pt, minimum height=0.95cm, text=white, font=\bfseries\fontsize{12}{13}\selectfont, inner sep=0pt},
  modulebox/.style={rounded corners=5pt, line width=0.9pt, fill=white, align=left, inner sep=7pt},
  handoff/.style={rounded corners=5pt, line width=0.9pt, align=center, inner sep=6pt, font=\fontsize{8.0}{9.1}\selectfont},
  flow/.style={-{Latex[length=2.8mm,width=2.2mm]}, draw=linegray, line width=1.05pt},
  vflow/.style={-{Latex[length=2.4mm,width=1.8mm]}, draw=muted!80, line width=0.9pt},
  note/.style={rounded corners=4pt, draw=linegray!80, fill=softgray, line width=0.7pt, align=left, inner sep=5pt, font=\fontsize{8.2}{9.2}\selectfont},
}

\newcommand{\modtitle}[1]{{\bfseries\fontsize{9.7}{11}\selectfont #1}}
\newcommand{\modbody}[1]{{\fontsize{8.45}{9.6}\selectfont #1}}
\newcommand{\bul}{\textbullet\;}



\DeclareLanguageMapping{american}{american-apa}
\addbibresource{materials_methods_piter_draft.bib}

\doublespacing
\setlength{\parindent}{0pt}
\setlength{\parskip}{0.5em}
\setlength{\tabcolsep}{5.4pt}
\renewcommand{\arraystretch}{1.42}
\setlength{\extrarowheight}{1.5pt}
\setlength{\emergencystretch}{2em}
\captionsetup{font={small,stretch=1},labelfont=bf}
\setlength{\arrayrulewidth}{1.0pt}
\arrayrulecolor{black}
\urlstyle{same}

\definecolor{stagecollect}{HTML}{365C8D}
\definecolor{stageclean}{HTML}{2E7D6E}
\definecolor{stageprep}{HTML}{C97A2B}
\definecolor{stagemine}{HTML}{7B4FA0}
\definecolor{lightcollect}{HTML}{EAF1F8}
\definecolor{lightclean}{HTML}{E8F4F1}
\definecolor{lightprep}{HTML}{FBF1E4}
\definecolor{lightmine}{HTML}{F1E9F8}
\definecolor{tableheader}{HTML}{F2F2F2}

\newcolumntype{L}[1]{>{\RaggedRight\arraybackslash}p{#1}}
\newcolumntype{Y}{>{\RaggedRight\arraybackslash}X}

\newcommand{\HlineFive}{\hhline{|-----|}}
\newcommand{\HlineSeven}{\hhline{|-------|}}
\newcommand{\cellstack}[1]{\shortstack[l]{#1}}

\newcommand{\stagecell}[2]{\cellcolor{#1!14}\textcolor{#1}{\textbf{#2}}}
\newcommand{\metriccell}[2]{\cellcolor{#1!14}\textbf{#2}}
\DeclareRobustCommand{\artifact}[1]{\texttt{#1}}
\newcommand{\cmd}[1]{\texttt{\detokenize{#1}}}

\title{Materials and Methods Draft\\Mouse DRG RNA-seq Follow-up Project}
\author{Piter Garcia}
\date{Draft prepared for BIOL550}

\begin{document}
\maketitle

\section{Introduction}

The published study behind this dataset investigates how injured neurons balance stress adaptation with regenerative growth, with particular emphasis on the aryl hydrocarbon receptor (AhR) as a regulator of that switch \parencite{halawani2023ahr}. In the source paper, AhR is presented as a neuronal brake on axon regeneration after injury, with mechanistic evidence linking AhR-dependent responses to stress preservation, proteostasis, and restraint of growth-associated programs \parencite{halawani2023ahr}. This biological context makes the dataset useful for reanalysis because it does not only capture whether injury changes gene expression, but also how those changes may relate to the larger transition between protective and regenerative states in sensory neurons \parencite{halawani2023ahr}.

The study focuses on dorsal root ganglion (DRG) neurons because they provide a strong model for peripheral axon regrowth after sciatic nerve injury, and the RNA-seq design compares ipsilateral DRGs from the injured side against contralateral DRGs from the uninjured side \parencite{halawani2023ahr}. The source paper also uses neuronal conditional knockout models of \textit{Ahr}, which are especially important here because they isolate neuron-intrinsic effects more cleanly than a constitutive knockout background with broader baseline pathology \parencite{halawani2023ahr}. In the reanalyzed subset used in this paper, genotype remains a relevant second layer of comparison between \artifact{ff} wild-type controls and \artifact{cre} Ahr conditional knockouts, but the side-specific injury contrast remains the main biological axis carried forward through the transcriptomic analysis \parencite{halawani2023ahr,sraStudySRP618841,bioprojectPRJNA1322439,geoGSE243308}.

Our paper uses the original study as a mechanistic anchor, but extends beyond its candidate-gene emphasis by applying transcriptome-wide differential expression, bend-point narrowing, and pathway-level analysis to the full RNA-seq dataset. Rather than repeating the paper's full experimental program, our aim is to test whether the strongest signals in the data remain centered on injury-side differences, evaluate how strongly genotype contributes relative to that main contrast, and connect the resulting gene sets to broader biological processes that support interpretation of the injury response \parencite{halawani2023ahr}. In that sense, this reanalysis uses the published dataset both to confirm the major biological hierarchy suggested by the source paper and to place the most reproducible expression changes into a broader pathway-level context.

\section{Materials and Methods}


This project reanalyzed mouse DRG bulk RNA-seq after sciatic nerve injury using a four-stage workflow modeled after a CRISP-style progression: \textbf{Data Collection}, \textbf{Data Cleaning}, \textbf{Data Preparation}, and \textbf{Data Analysis and Interpretation} \parencite{crispdm2000}. That structure reflects how the \artifact{mouse\_new} workflow was built and refined across the weekly deliverables, where each stage owned a specific deliverable, a checkpoint artifact, and a reusable output for the next stage. This staged structure made it possible to reuse the pipeline after earlier dataset-selection problems, because accession tracking, QC review, alignment review, and downstream modeling were separated into explicit, repeatable, and reproducible checkpoints.

The working subset comprised 20 NovaSeq X paired-end libraries from \artifact{SRP618841} (BioProject \artifact{PRJNA1322439}; GEO \artifact{GSE243308}), representing L4--L6 dorsal root ganglion samples collected 1 day after sciatic nerve injury \parencite{halawani2023ahr,sraStudySRP618841,bioprojectPRJNA1322439,geoGSE243308}. The local subset was organized as a balanced side-by-genotype design within the \artifact{family\_drg\_novaseqx} family, with equal representation of ipsilateral and contralateral samples and of \artifact{ff} (wild-type control) and \artifact{cre} (Ahr conditional knockout) genotype classes before quality control, alignment, and differential expression analysis. At each stage, outputs were reviewed and advanced through defined checkpoints to maintain consistency and reproducibility throughout the analysis.

\begin{figure}[H]
\centering
\makebox[\textwidth][c]{\includegraphics[width=1.02\textwidth]{assets_methods/BIOL550-Group_Project_drawio_light.png}}
\caption{Overview ribbon summarizing the four computational stages and their stage-owned handoffs in the \artifact{mouse\_new} workflow.}
\label{fig:overall_pipeline}
\end{figure}


\subsection{Data Cleaning: Quality Control and Read Trimming}
Raw read quality was evaluated with \artifact{FastQC} \parencite{fastqc}, and cohort-level summaries were consolidated with \artifact{MultiQC} \parencite{multiqc}. This stage was treated as a documented checkpoint, so pre- and post-trimming QC outputs were retained to show the effect of trimming explicitly. The local workflow used \artifact{fastp} v0.23.2 as the canonical cleanup tool because it provided a reproducible paired-end workflow, generated per-sample reports, and resolved adapter-related signal more effectively than the earlier baseline \parencite{fastp}.

The retained \artifact{fastp} configuration used paired-end adapter detection, poly-G trimming, a mean-quality cutoff of 20, a minimum retained length of 30 bases, and per-sample HTML/JSON reporting. In Figure~\ref{fig:data_cleaning}, the left panel shows the module-status heatmap and the right panel shows the SRR-level severity-change plot. Together, these artifacts document a consistent improvement in read quality across the cohort before alignment. Retaining both the raw and trimmed checkpoints also made this stage reusable for later dataset screening, because poor candidate subsets could be rejected or repaired using the same QC evidence.

\begin{figure}[H]
\centering
\makebox[\textwidth][c]{\includegraphics[width=1.04\textwidth]{assets_methods/data_cleaning_stage.png}}
\caption{Data Cleaning stage figure using retained QC artifacts and a rebuilt checkpoint strip.}
\label{fig:data_cleaning}
\end{figure}

\begin{table}[H]
\centering
\begin{singlespace}
\fontsize{7.2}{5.1}\selectfont
\caption{Data Cleaning stage table.}
\label{tab:data_cleaning}
\fontsize{6.9}{7.8}\selectfont
\begin{tabularx}{\textwidth}{|Y|Y|Y|Y|Y|}
\HlineFive
\rowcolor{tableheader}
\textbf{Substage} & \textbf{Input} & \textbf{Tool} & \textbf{Checkpoint artifact / evidence} & \textbf{Why it mattered} \\
\HlineFive
\stagecell{stageclean}{Raw QC} &
raw paired-end reads &
\artifact{FastQC} + \artifact{MultiQC} &
raw module-status heatmap and cohort QC summary; adapter-content failures across all 40 raw read files &
establishes the baseline quality state before trimming \\
\HlineFive

\stagecell{stageclean}{Read trimming} &
raw \artifact{FASTQ.gz} pairs &
\artifact{fastp} v0.23.2 &
\cellstack{trimmed read pairs\\per-sample HTML/JSON reports; PE adapter detect,\\Q20 cutoff, min length 30 bp,\\poly-G trimming} &
standardizes cleanup across all retained libraries \\
\HlineFive

\stagecell{stageclean}{Post-trim checkpoint} &
trimmed read pairs &
\artifact{FastQC} + \artifact{MultiQC} &
SRR-level severity summary: all retained severity deltas shifted in the negative direction after trimming &
justifies using the cleaned reads, not the raw reads, for alignment \\
\HlineFive
\end{tabularx}
\end{singlespace}
\end{table}

\subsection{Data Preparation: Reference Selection and Alignment}
Trimmed read pairs were aligned to the \textit{Mus musculus} \artifact{GRCm39} primary assembly with the matching Ensembl release 115 annotation \parencite{ensemblRelease115}. The reference pair used in the local workflow consisted of \artifact{Mus\_musculus.GRCm39.dna.primary\_assembly.fa} and \artifact{Mus\_musculus.GRCm39.115.gtf}. A STAR genome index was generated with \artifact{sjdbOverhang = 150} to match the 151-bp paired-end reads. Alignment was performed with STAR using gzip-aware input handling, coordinate-sorted BAM output, and GeneCounts, so that the same run also produced per-sample count tables \parencite{star}.

This stage linked cleaned reads to downstream modeling through a stable alignment handoff. For each sample, the workflow retained three key outputs: a coordinate-sorted BAM file, a STAR \artifact{Log.final.out} summary, and a \artifact{ReadsPerGene.out.tab} count table. These outputs provided both the alignment checkpoint and the count-level handoff used to assemble the final family matrix for the DRG subset. Reverse-stranded counts (column 4 of each \artifact{ReadsPerGene.out.tab} file) were used as the input for downstream DESeq2 modeling, consistent with the strandedness assumption used in the project workflow. Treating alignment review and count-matrix assembly as explicit handoffs also made the same preparation stage reusable across candidate datasets.

\begin{figure}[H]
\centering
\makebox[\textwidth][c]{\includegraphics[width=1.04\textwidth]{assets_methods/data_preparation_stage.png}}
\caption{Data Preparation stage figure built from retained STAR outputs and rebuilt alignment checkpoint metrics.}
\label{fig:data_preparation}
\end{figure}

\begin{table}[H]
\centering
\begin{singlespace}
\fontsize{7.2}{8.1}\selectfont
\caption{Data Preparation stage table.}
\label{tab:data_preparation}
\begin{tabularx}{\textwidth}{|Y|Y|Y|Y|Y|}
\HlineFive
\rowcolor{tableheader}
\textbf{Substage} & \textbf{Input} & \textbf{Tool} & \textbf{Checkpoint artifact / evidence} & \textbf{Why it mattered} \\
\HlineFive
\stagecell{stageprep}{Reference choice} &
mouse genome + annotation &
Ensembl reference selection &
reference FASTA/GTF pair for \artifact{GRCm39} primary assembly + Ensembl v115 annotation &
locks the coordinate system used by STAR and downstream counts \\
\HlineFive

\stagecell{stageprep}{STAR index \& run} &
trimmed read pairs + reference &
\artifact{STAR} &
sorted BAM, \artifact{Log.final.out}, and \artifact{ReadsPerGene.out.tab}; \artifact{sjdbOverhang = 150} and paired-end alignment/counting in one run &
produces the alignment checkpoint and the count-ready artifacts in one run \\
\HlineFive

\stagecell{stageprep}{Alignment checkpoint} &
saved STAR summaries &
local alignment review &
retained unique-mapping plot plus the STAR median summary table; unique 93.24\%, multi 5.10\%, \artifact{noFeature} 3.82\%, ambiguous 1.77\% &
confirms the cleaned reads mapped well enough to move into DESeq2 \\
\HlineFive

\stagecell{stageprep}{Count handoff} &
per-sample count outputs &
local merge scripts / notebooks &
family-level count matrix for \artifact{family\_drg\_novaseqx}; 20 samples retained in the final family matrix &
converts per-sample STAR outputs into the DE-ready family matrix \\
\HlineFive
\end{tabularx}
\end{singlespace}
\end{table}

\subsection{Data Analysis and Interpretation: Differential Expression and Functional Analysis}
The merged family count matrix was analyzed with DESeq2 using the explicit side-by-genotype design carried forward from the sample manifest \parencite{deseq2}. An interaction design formula (\artifact{\textasciitilde{} side\_class * geno\_class}) was specified so that the effect of side could differ by genotype, while all five contrast branches were extracted from a single fitted model. Before modeling, low-support genes were filtered, reducing the tested set from 78{,}334 to 21{,}481 genes while retaining all 20 DRG samples. Five contrasts were then evaluated: \artifact{ipsi\_vs\_contra\_in\_ff}, \artifact{ipsi\_vs\_contra\_in\_cre}, \artifact{geno\_in\_contra}, \artifact{geno\_in\_ipsi}, and \artifact{interaction}. This contrast inventory was kept fixed so that later interpretation did not drift when one branch became more interpretable than another.

Downstream interpretation followed a PCA-first, branch-priority logic. PCA was used to evaluate sample structure before long DE tables were emphasized, and the side-specific branches were then carried forward as the primary follow-up path. To avoid reducing very large DE outputs to unstructured gene lists, the workflow applied a bend-point rule to the ordered adjusted-\textit{p}-value curve for each contrast. These bend-point-derived follow-up sets were reported alongside the \artifact{padj < 0.05} results. For the two primary side-specific contrasts, this step retained 709 and 870 genes, respectively.

The resulting follow-up sets were then submitted to g:Profiler using GO Biological Process, KEGG, and Reactome layers \parencite{gprofiler2021,geneOntology2021,kegg2025,reactome2022}. In the \artifact{ff} branch, which remained the clearest pathway-level follow-up, enrichment outputs were further summarized using shared-gene chart, tree, and network views to reduce redundancy among overlapping GO labels. These enrichment products are described here as downstream summaries of the differential-expression workflow rather than as final biological claims.

\begin{figure}[H]
\centering
\makebox[\textwidth][c]{\includegraphics[width=1.04\textwidth]{assets_methods/data_mining_stage.png}}
\caption{Rebuilt Data Analysis and Interpretation stage figure derived from saved \artifact{DESeq2}, bend-point, and \artifact{g:Profiler} summary tables.}
\label{fig:data_analysis}
\end{figure}

\begin{figure}[H]
\centering
\makebox[\textwidth][c]{\includegraphics[width=1.04\textwidth]{assets_methods/data_mining_selection_stage.png}}
\caption{Retained bend-point checkpoint summary for the two primary side-specific branches (cards). The ordered p-value and cumulative-count plots are shown for the \artifact{ipsi\_vs\_contra\_in\_ff} example, with the right panel zoomed to \textit{p} $\le$ 0.01.}
\label{fig:data_analysis_selection}
\end{figure}

\begin{table}[H]
\centering
\begin{singlespace}
\fontsize{7.1}{8.0}\selectfont
\caption{Data Analysis and Interpretation workflow table.}
\label{tab:data_analysis}
\begin{tabularx}{\textwidth}{|Y|Y|Y|Y|Y|}
\HlineFive
\rowcolor{tableheader}
\textbf{Substage} & \textbf{Input} & \textbf{Tool} & \textbf{Checkpoint artifact / evidence} & \textbf{Why it mattered} \\
\HlineFive
\stagecell{stagemine}{Model setup} &
family count matrix + design table &
\artifact{DESeq2} &
normalized counts, VST matrix, and contrast-level result tables; 20 samples retained and 78{,}334 $\rightarrow$ 21{,}481 genes after filtering &
defines the inferential base for every contrast and downstream plot \\
\HlineFive

\stagecell{stagemine}{Contrast production} &
filtered model outputs &
\cellstack{\artifact{DESeq2}\\contrast manifest} &
five saved contrast branches; two primary side-specific and three supporting genotype / interaction branches &
keeps the main and supporting stories explicit before list narrowing \\
\HlineFive

\stagecell{stagemine}{Bend-point selection} &
ordered adjusted-\textit{p} tables &
local bend-point scripts &
analysis-summary and per-contrast bend-point summary tables; primary side-specific follow-up sets of 709 genes in \artifact{ff} and 870 genes in \artifact{cre} &
turns very large DE lists into smaller reproducible follow-up sets \\
\HlineFive

\stagecell{stagemine}{Enrichment layer} &
bend-point-selected gene sets &
\artifact{g:Profiler} &
g:Profiler source-summary and enrichment tables; the main \artifact{ff} branch retained 359 GO:BP, 18 KEGG, and 3 Reactome terms &
captures the terminal functional-summary products without replacing the gene-level story \\
\HlineFive
\end{tabularx}
\end{singlespace}
\end{table}

\begin{table}[H]
\centering
\begin{singlespace}
\fontsize{7.2}{8.0}\selectfont
\caption{Primary side-specific follow-up summary carried out during the Data Analysis stage.}
\label{tab:data_analysis_primary}
\begin{tabularx}{\textwidth}{|Y|Y|Y|Y|Y|Y|Y|}
\HlineSeven
\rowcolor{tableheader}
\textbf{Branch} & \textbf{\artifact{padj < 0.05} genes} & \textbf{Bend-point genes} & \textbf{Threshold} & \textbf{GO:BP} & \textbf{KEGG} & \textbf{REAC} \\
\HlineSeven
\stagecell{stagemine}{ff side branch} &
\metriccell{stagemine}{7{,}023} &
\metriccell{stagemine}{709} &
\artifact{1.37e-17} &
359 &
18 &
3 \\
\HlineSeven

\stagecell{stagemine}{cre side branch} &
\metriccell{stagemine}{7{,}541} &
\metriccell{stagemine}{870} &
\artifact{8.40e-17} &
425 &
21 &
7 \\
\HlineSeven
\end{tabularx}
\end{singlespace}
\end{table}

\section{Results}

\subsection{Quality control and trimming}

Quality control was reviewed first with paired raw-versus-trimmed \artifact{FastQC} and \artifact{MultiQC} outputs before any downstream modeling was interpreted. The strongest and most consistent raw-data problem was adapter-associated signal, which was expected for this RNA-seq workflow but still needed to be reduced before alignment. In the saved notebook outputs, the pre-trim cohort showed widespread adapter-content failures, while the trimmed cohort showed a sharp reduction in adapter burden after the \artifact{fastp} cleanup step. This made the cleaned read set acceptable for downstream alignment while still preserving the broader RNA-seq context that duplication and some overrepresented-sequence signals were not expected to disappear completely.

\begin{figure}[H]
\centering
\makebox[\textwidth][c]{\includegraphics[width=0.90\textwidth]{assets_htsa/qc_adapter_pre_post.png}}
\caption{Adapter-content comparison before and after trimming for the retained \artifact{mouse\_new} cohort. The pre-trim panel shows the widespread adapter signal that motivated cleanup, and the post-trim panel shows the reduced adapter burden carried into alignment.}
\label{fig:results_qc_adapter}
\end{figure}

\subsection{Alignment quality and global sample structure}

Once the reads were cleaned, alignment to the \textit{Mus musculus} GRCm39 reference showed a consistently high-quality mapping profile across the retained 20-sample subset. The unique-mapping percentage stayed above 90\% for every sample, with a median unique-mapping rate of 93.24\% and a narrow interquartile range of 93.10--93.42\%. This level of consistency supported keeping the subset for transcriptome-wide modeling rather than dropping runs for technical instability. The same conclusion carried into the sample-structure review, where the PCA and sample-to-sample distance heatmap both showed that the dominant visible organization in the dataset was by injury side, while genotype appeared as a weaker secondary layer.

\begin{figure}[H]
\centering
\begin{minipage}[t]{0.49\textwidth}
\centering
\includegraphics[width=\textwidth]{assets_htsa/alignment_unique_mapping_bar.png}
\end{minipage}\hfill
\begin{minipage}[t]{0.49\textwidth}
\centering
\includegraphics[width=\textwidth]{assets_htsa/alignment_unique_mapping_box.png}
\end{minipage}
\caption{Alignment-quality summary for the retained \artifact{mouse\_new} subset. Left: per-sample unique-mapping percentages. Right: distribution summary of unique-mapping rates across the cohort.}
\label{fig:results_alignment_mapping}
\end{figure}

\begin{figure}[H]
\centering
\begin{minipage}[t]{0.49\textwidth}
\centering
\includegraphics[width=\textwidth]{assets_htsa/results_pca.png}
\end{minipage}\hfill
\begin{minipage}[t]{0.49\textwidth}
\centering
\includegraphics[width=\textwidth]{assets_htsa/results_sample_distance_heatmap.png}
\end{minipage}
\caption{Global structure of the retained DRG RNA-seq subset. Left: PCA of the 20-sample cohort. Right: sample-to-sample distance heatmap from transformed count data. Together these figures support the same hierarchy: side is the clearest organizing variable, while genotype is present but secondary.}
\label{fig:results_pca_heatmap}
\end{figure}

\subsection{Differential expression and bend-point narrowing}

Differential expression was modeled in \artifact{DESeq2} using the design \artifact{\textasciitilde{} genotype + side + genotype:side}, and the strongest output again came from the two side-specific contrasts. To avoid treating the entire \artifact{padj < 0.05} output as one equally informative block, the notebook pipeline then ranked genes by adjusted \textit{p}-value and used the ordered adjusted-\textit{p}-value curve as a reproducible bend-point rule. For \artifact{ipsi\_vs\_contra\_in\_ff}, that reduced 7{,}023 significant genes to a 709-gene follow-up set; for \artifact{ipsi\_vs\_contra\_in\_cre}, it reduced 7{,}541 genes to 870. This narrowing step did not remove the underlying biological signal. Instead, it defined smaller contrast-specific cores that were still strongly supported statistically and were practical for pathway interpretation.

\begin{figure}[H]
\centering
\makebox[\textwidth][c]{\includegraphics[width=0.92\textwidth]{assets_htsa/results_bendpoint_volcano.png}}
\caption{Bend-point-guided narrowing for the main WT branch, \artifact{ipsi\_vs\_contra\_in\_ff}. The comparison figure contrasts the full significant set with the narrower bend-point-selected core used for downstream interpretation.}
\label{fig:results_bendpoint}
\end{figure}

\subsection{Gene Ontology and pathway-level interpretation}

The clearest enrichment result came from the WT-like branch, \artifact{ipsi\_vs\_contra\_in\_ff}, especially after the upregulated portion of the bend-point-selected set was examined with retained term-to-gene membership. The strongest recurring Biological Process themes centered on signaling, migration, morphogenesis, response to stimulus, and neurogenic remodeling. These themes are visible in the composite GO summary and become easier to interpret once the overlap among terms is shown directly rather than treated as a flat list of independent labels.

\begin{figure}[H]
\centering
\makebox[\textwidth][c]{\includegraphics[width=0.92\textwidth]{assets_htsa/results_go_composite.png}}
\caption{GO summary derived from the bend-point-selected WT-side branch. The retained views highlight the strongest enriched terms and the overlap structure that keeps the pathway story compact enough to explain in prose.}
\label{fig:results_go_composite}
\end{figure}

\begin{figure}[H]
\centering
\begin{minipage}[t]{0.49\textwidth}
\centering
\includegraphics[width=\textwidth]{../differential_expression_all20/derived_analysis/ipsi_vs_contra_in_ff/shinygo_style/ff_up_network_view.png}

\vspace{0.3em}
\fontsize{9}{10}\selectfont WT branch network view (\artifact{ipsi\_vs\_contra\_in\_ff})
\end{minipage}\hfill
\begin{minipage}[t]{0.49\textwidth}
\centering
\includegraphics[width=\textwidth]{../differential_expression_all20/derived_analysis/ipsi_vs_contra_in_cre/shinygo_style_shared/cre_up_network_view.png}

\vspace{0.3em}
\fontsize{9}{10}\selectfont cKO branch network view (\artifact{ipsi\_vs\_contra\_in\_cre})
\end{minipage}
\caption{Shared-gene GO network views for the two main side-specific branches. Showing both networks side by side makes it easier to compare the cleaner WT-centered pathway story against the slightly broader but still informative cKO-side branch.}
\label{fig:results_go_networks}
\end{figure}

\textbf{Methods used for the GO follow-up.} Following the previous report and the later transcript guidance to keep the story centered on the modeled side-specific contrasts, the pathway analysis began from the bend-point-selected gene sets already established in the \artifact{DESeq2} notebook rather than from a new ad hoc comparison. The main differential-expression pipeline was run in \artifact{mouse\_differential\_expression\_all20.ipynb}, where the \artifact{DESeq2} design, PCA review, volcano outputs, and ordered adjusted-\textit{p}-value bend-point thresholds were all generated from the same retained 20-sample DRG subset. The GO follow-up was then rebuilt in \artifact{mouse\_differential\_expression\_ff\_shinygo\_style.ipynb} by re-running GO Biological Process enrichment through the \artifact{g:Profiler} API with retained term-to-gene membership and generating local chart, tree, network, groups, and gene-membership views. Consistent with the transcript emphasis that less is more and that retained figures must remain explainable to a non-specialist reader, only the smallest set of GO visuals that materially improved interpretation was carried forward into the paper.

\textbf{WT versus cKO branch comparison.} Among the follow-up branches, the WT-like side-specific branch \artifact{ipsi\_vs\_contra\_in\_ff} remained the clearest pathway-level result. It combined a strong side-driven separation, a large but reducible DE set, and the most coherent enrichment profile. GO Biological Process dominated the enrichment summary, and the later FF GO follow-up clarified that the apparent abundance of terms was not best understood as many unrelated pathways. Instead, once shared-gene structure was examined through the chart, tree, and network views, the strongest upregulated terms clustered around signaling, migration, morphogenesis, stress-response, and neurogenic remodeling themes. This convergence made the FF branch the most defensible biological centerpiece for the paper’s main interpretation. By contrast, the cKO-like branch \artifact{ipsi\_vs\_contra\_in\_cre} still supports the same side-specific biological logic but is slightly less clean as a headline pathway result, while \artifact{geno\_in\_contra} works best as the strongest secondary genotype comparison rather than as a co-equal story.

\section{Discussion}

\subsection{Main biological interpretation}

\textbf{Overall pattern in the data.} Across the weekly reports, the most stable shared finding was that the side-specific injury contrasts carried the clearest structure in the \artifact{mouse\_new} dataset. The early alignment and QC checkpoints supported using this 20-sample DRG subset for downstream modeling, but the stronger discussion point emerged only after the PCA-first review: side separated more clearly than genotype, while genotype behaved as a secondary axis with substantial overlap between paired \artifact{ff}/\artifact{cre} samples. That pattern was already visible in the project’s first differential-expression checkpoint and remained consistent through the later bend-point and GO follow-up reports. Nikhi’s intermediate report reached the same practical conclusion from the PCA and sample-distance views, namely that ipsilateral-versus-contralateral structure was clearer than genotype structure and therefore more defensible as the main biological path for interpretation.

\begin{figure}[H]
\centering
\makebox[\textwidth][c]{\includegraphics[width=0.92\textwidth]{assets_discussion/pca_side_genotype_annotated.png}}
\caption{PCA-first view of the \artifact{family\_drg\_novaseqx} subset used in Discussion. The dominant separation occurs along the side axis rather than by genotype, supporting the decision to prioritize side-specific branches in the downstream interpretation.}
\label{fig:discussion_pca}
\end{figure}

\textbf{Strongest supported signal; primary versus secondary findings.} The PCA-first reading mattered because it disciplined how the large DE outputs were interpreted. Instead of treating every modeled contrast as equally informative, the reports converged on a branch-priority logic in which \artifact{ipsi\_vs\_contra\_in\_ff} and \artifact{ipsi\_vs\_contra\_in\_cre} served as the primary follow-up branches, while \artifact{geno\_in\_contra} was retained as the best supporting genotype comparison. This common structure also explains why the discussion should emphasize what the reports share instead of over-weighting report-specific exploratory details. The secondary branches are still informative, especially when they help show that genotype effects are not absent, but they do not displace the side-specific injury signal as the central story.

\subsection{Gene-expression differences across the main biological contrasts}

\textbf{Ipsilateral versus contralateral.} The bend-point step was the key transition from broad differential-expression output to a more interpretable follow-up set. Both the weekly reports and the internal notebooks used the ordered adjusted-\textit{p}-value curve as a reproducible narrowing rule rather than relying only on arbitrary top-\(N\) lists or on raw \artifact{padj < 0.05} counts. For the two main side-specific contrasts, that rule reduced 7{,}023 and 7{,}541 significant genes to 709 and 870 genes, respectively, producing smaller sets that were still strongly supported statistically and were practical for pathway analysis. Nikhi’s follow-up report highlighted the same value of this reduction step: it preserved a highly significant subset while making pathway-level interpretation manageable enough to compare against known injury and regeneration mechanisms.

\begin{figure}[H]
\centering
\makebox[\textwidth][c]{\includegraphics[width=0.92\textwidth]{assets_discussion/ordered_pvalue_and_cumulative_curve.png}}
\caption{Bend-point diagnostic for \artifact{ipsi\_vs\_contra\_in\_ff}. The ordered adjusted-\textit{p}-value curve and cumulative-count view show how the follow-up set was reduced to a reproducible core rather than chosen as an arbitrary top-gene list.}
\label{fig:discussion_bendpoint}
\end{figure}

\begin{figure}[H]
\centering
\makebox[\textwidth][c]{\includegraphics[width=0.90\textwidth]{assets_discussion/ff_up_tree_view.png}}
\caption{Shared-gene tree view for the upregulated \artifact{ipsi\_vs\_contra\_in\_ff} GO follow-up. The clustering pattern supports interpreting the FF enrichment result as a connected injury-response neighborhood rather than as many independent pathway stories.}
\label{fig:discussion_ff_tree}
\end{figure}

\subsection{Biological processes implicated by the analysis}

The most defensible process-level interpretation is that the dataset reflects a broad injury-response transcriptional program rather than a small number of isolated candidate-gene effects. The strongest shared terms across the FF follow-up involve signaling, migration, morphogenesis, response to stimulus, and neurogenic or developmental remodeling. These themes fit the idea that the injured DRG state is being captured as a coordinated systems-level response after nerve injury rather than as a single-pathway perturbation.

Supporting terms in the saved enrichment tables suggest additional process classes that are worth discussing more cautiously. Oxidative-stress and related response terms appear in both the FF upregulated follow-up and in supporting genotype-related branches, which makes stress-response and proteostasis-adjacent regulation reasonable secondary themes. Translation-related terms are also present, especially translational initiation and regulation-of-translation signals in the bend-point-selected outputs, suggesting that injury-associated remodeling may include active control over protein-synthesis programs rather than only transcription-factor-level changes. Metabolic and energy-related terms are weaker and less central in the main FF story, but they appear in supporting branches and therefore may still help contextualize how the transcriptomic response coordinates signaling with cellular resource use. These process classes should be treated as transcriptomic associations, not as direct proof of mechanism, and any AhR-related or other pathway-specific hypotheses should remain tentative unless they are explicitly supported by the selected gene memberships and by later biological validation.

\subsection{What was unusual or needs caution}

Several patterns still require caution. First, some volcano plots are visibly asymmetric, which supports the idea of directionally biased injury-response changes but also warns against treating the entire significant set as one uniform biological block. Second, the genotype signal appears weaker and more secondary than the side signal, a conclusion supported by the PCA overlap and the near-collision pairs discussed in the reports. Third, the visible sample structure can support interpretation, but it does not by itself establish causation. The current pipeline therefore supports prioritization and biological framing, not mechanistic proof.

\subsection{What this adds beyond the original paper}

This reanalysis adds value beyond the original paper by expanding the interpretation from highlighted candidate genes to a broader transcriptomic context. In practice, that means the workflow does three useful things: it extends the analysis beyond a small candidate list, it links pathways that would otherwise appear as disconnected labels, and it connects newly identified genes back to the genes emphasized in the original study through shared pathway and network structure. In that sense, the NGS-based analysis contributes both global discovery and reproducibility. It does not replace the paper’s focal genes or experimental logic, but it helps show how those genes sit inside a wider injury-response program and which additional genes or pathways may be biologically linked to them.

\subsection{Future validation and limitations}

The strongest immediate follow-up would be targeted validation of key differentially expressed genes by qPCR, especially for genes that anchor the main side-specific branches and for genes that help connect the original paper’s candidates to the broader pathway-level result. Functional follow-up experiments would then be the next step for the most plausible injury-response candidates. Two-hybrid screening or similar interaction-focused assays should only be pursued if the selected genes and pathway context make that biologically justified, rather than being treated as a default next step.

The main limitations remain the same across the reports: the genotype signal is weaker than the side signal, PCA overlap limits how strongly genotype-specific effects can be claimed, the analysis depends on the original dataset design, and pathway interpretation should not be confused with causal proof. GO redundancy is especially important here, because overlapping annotations can exaggerate how many distinct biological stories are present unless the shared-gene structure is examined directly. Within those limits, however, the weekly reports support a coherent conclusion: the \artifact{mouse\_new} workflow converged on a robust side-specific injury-response program, and the FF branch provided the clearest pathway-level window into that response.

\printbibliography[title={References}]

\end{document}
